% ============================================================================
% Fichier  : partie3/P03_C09_referentiels_correspondances.tex
% Titre    : Les Autres Référentiels et leurs Correspondances
% Auteur   : MINKA MI NGUIDJOÏ Thierry Emmanuel
% Version  : 1.0 -- novembre 2026
% Corrigés : corriges/Corr_C09.tex
% ============================================================================
% Sources : NIST CSF 2.0 (2024 : 6 fonctions, 22 catégories, 106
% sous-catégories, quatre niveaux de mise en œuvre, profils) ; NIST SP
% 800-53 rév. 5 (20 familles) ; CIS Controls v8 (2021, 18 contrôles, 153
% mesures, 3 groupes de mise en œuvre, IG1 = 56 mesures ; révision 8.1 de
% 2024) ; PCI DSS v4.0 (2022) ; critères de services de confiance de l'AICPA
% (SOC 2) ; ISAE 3402 (types 1 et 2) ; COBIT 2019 (40 objectifs). Chiffres
% d'après les publications des éditeurs, à confirmer par l'auteur ;
% intitulés traduits par le cours.
% Faits du cas (bible et corpus) : ANPP (Yaoundé), agence publique de
% paiement contrôlée par la Chambre des comptes ; M. ESSOMBA, directeur
% général ; M. NJOYA, magistrat vérificateur ; 4 agents payés après leur
% départ ; aucune procédure écrite de désactivation des comptes ; pièces
% « aux archives ».
% Faits nouveaux (reportés dans la bible) : en 2027, un bailleur de fonds
% exige un profil NIST CSF 2.0 ; base commune de 10 contrôles B1 à B10 ;
% résultats 4/3/3 ; projection par fonction du CSF ; niveau actuel 1
% « partiel », cible 3 ; DIGICORE pourrait remettre un rapport ISAE 3402 de
% type 2.
% ============================================================================

\chapter{Les Autres Référentiels et leurs Correspondances}
\label{chap:referentiels}
\chaptermark{Référentiels et correspondances}

\epigraph{%
    « Tester une fois,\\
    rendre compte autant de fois qu'il le faut. »%
}{Principe du cours}

\niveauChapitre{Fondamental}{%
    Maîtrise du Chapitre~\ref{chap:iso27000}%
}

\begin{boxobjectifs}
\begin{itemize}
    \item Présenter le cadre de cybersécurité du NIST (CSF~2.0), le
          catalogue NIST SP~800-53, les CIS Controls et PCI DSS.
    \item Situer les rapports d'assurance SOC~2 et ISAE~3402, et COBIT~2019.
    \item Construire une table de correspondance entre référentiels, et en
          connaître les pièges.
    \item Bâtir une base commune de contrôles, testée une fois et projetée
          sur plusieurs référentiels.
    \item Choisir les référentiels d'une organisation.
\end{itemize}
\end{boxobjectifs}

\bigskip

% ============================================================================
\section*{Avant de commencer : trois audits pour un seul système}
% ============================================================================

Yaoundé, mars 2027. L'ANPP\footnote{L'ANPP, ses dirigeants, le bailleur,
les dates et les montants sont fictifs~; les référentiels cités sont
exacts.}, agence publique de paiement, finance la modernisation de ses
systèmes avec l'appui d'un bailleur de fonds international. La convention
de financement exige un « profil de cybersécurité selon le cadre du NIST,
version~2.0 », actuel et cible. La Chambre des comptes, par la voix du
magistrat vérificateur M.~NJOYA, prépare son contrôle des paiements et de
leurs systèmes. Et l'audit de sécurité prévu par la loi n°~2010/012 pourrait
s'appliquer à l'agence (Chapitre~\ref{chap:droit}).

M.~ESSOMBA, directeur général, refuse d'organiser trois chantiers~: «~Nos
serveurs sont les mêmes pour tout le monde. Je veux un seul travail, et
trois rapports.~»

\medskip

Ce chapitre présente les principaux référentiels qui coexistent avec
l'ISO/IEC~27001, puis la technique qui permet d'en satisfaire plusieurs à
la fois~: la table de correspondance et la base commune de contrôles.

\begin{boxalert}
\textbf{Violation typique.} Conclure qu'une organisation est conforme au
cadre du NIST parce qu'elle est certifiée ISO/IEC~27001, au motif qu'une
table de correspondance relie les deux.

\smallskip
\textbf{Règle~:} une correspondance dit que deux exigences portent sur le
même sujet, pas qu'elles disent la même chose. La conformité se démontre
référentiel par référentiel, à partir des preuves.
\end{boxalert}

\separateur

% ============================================================================
\section{Le Cadre de Cybersécurité du NIST}
\label{sec:c09-csf}
% ============================================================================

\subsection{Six fonctions, vingt-deux catégories}

Le cadre de cybersécurité du NIST (\sigle{CSF}), dans sa version~2.0 de
2024, organise la gestion des risques cyber en six fonctions, déclinées en
22 catégories et 106 sous-catégories~\cite{nist2024csf}. Le
Chapitre~\ref{chap:grc} a présenté la fonction \emph{Gouverner}~; voici
l'ensemble.

\begin{tableaucours}{L{2.4cm}L{9.8cm}}{%
    \tableauHeader{Fonction} &
    \tableauHeader{Catégories (intitulés traduits par le cours)}
}
Gouverner (GV) & Contexte organisationnel (OC)~; stratégie de gestion des
risques (RM)~; rôles, responsabilités et autorités (RR)~; politique (PO)~;
supervision (OV)~; risques de la chaîne d'approvisionnement (SC) \\
Identifier (ID) & Gestion des actifs (AM)~; appréciation des risques
(RA)~; amélioration (IM) \\
Protéger (PR) & Identités, authentification et contrôle d'accès (AA)~;
sensibilisation et formation (AT)~; sécurité des données (DS)~; sécurité
des plateformes (PS)~; résilience de l'infrastructure (IR) \\
Détecter (DE) & Surveillance continue (CM)~; analyse des événements
indésirables (AE) \\
Répondre (RS) & Gestion des incidents (MA)~; analyse (AN)~; communication
(CO)~; atténuation (MI) \\
Rétablir (RC) & Exécution du plan de rétablissement (RP)~; communication
(CO) \\
\end{tableaucours}
\vspace{-0.8em}
\begin{center}{\small\itshape Tableau~9.1. Les fonctions et catégories du
NIST CSF~2.0}\end{center}

\subsection{Profils et niveaux}

Le CSF ne fixe pas d'exigences~; il donne un langage. Une organisation
décrit son \emph{profil actuel} (ce qu'elle fait, sous-catégorie par
sous-catégorie) et son \emph{profil cible} (ce qu'elle vise), puis comble
l'écart. Elle peut aussi caractériser la rigueur de sa gestion des risques
par l'un des quatre \emph{niveaux de mise en œuvre}~: partiel, informé par
le risque, reproductible, adaptatif.

\begin{boxinfo}
\textbf{Pourquoi les bailleurs aiment le CSF.} Il est gratuit, lisible par
une direction, et ne suppose pas de certification. Un profil actuel et un
profil cible donnent au bailleur une feuille de route qu'il peut suivre
d'une année sur l'autre. C'est aussi sa limite~: un profil est une
auto-description~; il ne vaut que par les preuves qui l'appuient.
\end{boxinfo}

% ============================================================================
\section{Catalogues de Mesures}
\label{sec:c09-catalogues}
% ============================================================================

\subsection{NIST SP~800-53}

Le catalogue NIST SP~800-53, révision~5, compte plus d'un millier de
mesures et d'améliorations, réparties en vingt familles (contrôle d'accès,
audit et responsabilité, réponse aux incidents, etc.)~\cite{nist2020sp80053}.
Il est obligatoire pour les systèmes fédéraux américains et sert ailleurs
de réservoir de mesures détaillées. Son compagnon, le SP~800-53A, décrit la
manière d'évaluer chaque mesure (Chapitre~\ref{chap:audit}).

\subsection{Les CIS Controls}

Les CIS Controls, version~8, publiés par le Center for Internet Security,
retiennent 18 contrôles et 153 mesures, classés par ordre de
priorité~\cite{cis2021controls}. Leur originalité est de définir trois
\emph{groupes de mise en œuvre}~: le premier (IG1) réunit 56 mesures,
présentées comme l'hygiène de base que toute organisation devrait appliquer~;
les groupes~2 et~3 ajoutent des mesures pour les organisations plus
exposées ou plus outillées.

\begin{tableaucours}{cL{5.0cm}cL{5.0cm}}{%
    \tableauHeader{N°} &
    \tableauHeader{Contrôle (traduit)} &
    \tableauHeader{N°} &
    \tableauHeader{Contrôle (traduit)}
}
1 & Inventaire des actifs matériels & 10 & Défenses contre les maliciels \\
2 & Inventaire des logiciels & 11 & Restauration des données \\
3 & Protection des données & 12 & Gestion de l'infrastructure réseau \\
4 & Configuration sécurisée & 13 & Surveillance et défense du réseau \\
5 & Gestion des comptes & 14 & Sensibilisation et formation \\
6 & Gestion du contrôle d'accès & 15 & Gestion des prestataires \\
7 & Gestion continue des vulnérabilités & 16 & Sécurité des applications \\
8 & Gestion des journaux & 17 & Réponse aux incidents \\
9 & Protection de la messagerie et du navigateur & 18 & Tests d'intrusion \\
\end{tableaucours}
\vspace{-0.8em}
\begin{center}{\small\itshape Tableau~9.2. Les 18 CIS Controls (version~8)}\end{center}

\begin{boxinfo}
\textbf{Un point de départ pour les petites organisations.} Pour une
entreprise comme LiZbiZ, qui n'a ni RSSI ni audit interne, le groupe IG1
des CIS Controls est un objectif réaliste et vérifiable avant toute
certification. Le Chapitre~\ref{chap:audit-technique} s'appuie sur les
guides de configuration du même organisme.
\end{boxinfo}

\subsection{PCI DSS}

La norme PCI DSS, version~4.0, est imposée par les réseaux de cartes de
paiement à toute organisation qui stocke, traite ou transmet des données
de cartes~\cite{pcissc2022dss}. Elle est contractuelle, détaillée et
vérifiée par des évaluateurs agréés. Elle ne concerne pas l'ANPP, qui paie
par virement~; elle concerne la BANQUE SANAGA pour son activité monétique,
sur un périmètre (l'environnement des données de cartes) distinct de celui
de son SMSI.

% ============================================================================
\section{Rapports d'Assurance et Gouvernance}
\label{sec:c09-assurance}
% ============================================================================

\subsection{SOC~2 et ISAE~3402}

Quand une organisation confie une activité à un prestataire, elle ne peut
pas toujours l'auditer elle-même. Elle peut alors demander un
\termecle{rapport d'assurance}\index{rapport d'assurance} établi par un
auditeur indépendant sur les contrôles du prestataire.

\begin{tableaucours}{L{2.4cm}L{4.6cm}L{5.2cm}}{%
    \tableauHeader{Rapport} &
    \tableauHeader{Cadre} &
    \tableauHeader{Contenu}
}
ISAE~3402 & Norme internationale d'assurance~\cite{iaasb2011isae3402} &
Contrôles d'un prestataire qui affectent l'information financière de ses
clients \\
SOC~2 & Critères de services de confiance de l'AICPA~\cite{aicpa2017tsc} &
Sécurité, disponibilité, intégrité des traitements, confidentialité, vie
privée \\
\end{tableaucours}
\vspace{-0.8em}
\begin{center}{\small\itshape Tableau~9.3. Deux rapports d'assurance sur les
prestataires}\end{center}

Chacun existe en deux types~: le \emph{type~1} porte sur la conception des
contrôles à une date~; le \emph{type~2} sur leur efficacité pendant une
période, généralement de six à douze mois. Seul le type~2 dit si les
contrôles ont fonctionné.

\begin{boxscenario}{Opération SANAGA. Une issue pour DIGICORE}
\textbf{Faits.} DIGICORE refuse l'avenant qui donnerait à la banque un droit
d'audit (Chapitre~\ref{chap:iso27005}).

\medskip
\textbf{Analyse} (confiance moyenne). Un prestataire qui sert plusieurs
banques refuse souvent d'être audité par chacune. La banque peut exiger, à
la place, un rapport d'assurance de type~2 par an, établi par un auditeur
indépendant, sur les contrôles qui la concernent. La base \NEK{} de
R-DC-02 deviendrait alors une base \OUI{} ou \NON{}, sans que la banque
entre chez le prestataire.
\end{boxscenario}

\subsection{COBIT~2019}

COBIT~2019, de l'ISACA, est un cadre de gouvernance et de gestion de
l'informatique, qui compte quarante objectifs~\cite{isaca2018cobitintro}.
Il n'est pas spécifique à la sécurité~: il la situe dans l'ensemble du
pilotage du système d'information. Le cours GCI (MSN-2S3-MSI03) l'étudie en
détail~; NRAC n'en retient que la distinction entre gouvernance et gestion
(Chapitre~\ref{chap:grc}).

% ============================================================================
\section{Les Tables de Correspondance}
\label{sec:c09-correspondances}
% ============================================================================

\subsection{Ce qu'est une correspondance}

\begin{boxdef}
\termecle{Table de correspondance}\index{table de correspondance}~:
tableau qui relie les exigences d'un référentiel à celles d'un autre, en
indiquant, pour chaque lien, s'il est complet ou partiel. Elle sert à
réutiliser un travail fait pour un référentiel au profit d'un autre.
\end{boxdef}

\begin{boxprocedure}
\textbf{Construire une correspondance fiable.} Le papier de travail qui
établit ou utilise une table de correspondance vérifie~:
\begin{itemize}
    \item la \textbf{source} de chaque lien~: publication officielle (le
          NIST et le CIS publient des correspondances), travail interne,
          ou fournisseur~;
    \item la \textbf{nature} du lien~: complet (les deux exigences ont le
          même contenu) ou partiel (l'une couvre une partie de l'autre)~;
    \item le \textbf{sens} du lien~: satisfaire A peut suffire pour B sans
          que l'inverse soit vrai~;
    \item la \textbf{multiplicité}~: une exigence peut correspondre à
          plusieurs autres, et réciproquement~;
    \item la \textbf{version} de chaque référentiel~: une table ISO/IEC
          27002:2013 ne vaut pas pour l'édition de 2022.
\end{itemize}
\end{boxprocedure}

\subsection{La base commune de l'ANPP}

Le responsable de la sécurité de l'ANPP construit une \termecle{base commune
de contrôles}\index{base commune de contrôles}~: une liste de contrôles
propres à l'agence, chacun relié aux référentiels qui l'exigent.

\begin{tableaucours}{cL{4.2cm}ccc}{%
    \tableauHeader{Réf.} &
    \tableauHeader{Contrôle de l'ANPP} &
    \tableauHeader{ISO/IEC 27002} &
    \tableauHeader{CSF~2.0} &
    \tableauHeader{CIS v8}
}
B1 & Inventaire des postes et serveurs & 5.9 & ID.AM & 1 \\
B2 & Désactivation des comptes au départ & 5.18 & PR.AA & 5 \\
B3 & Authentification multifacteur & 8.5 & PR.AA & 6 \\
B4 & Journalisation et revue & 8.15 & DE.CM & 8 \\
B5 & Sauvegardes testées & 8.13 & PR.DS, RC.RP & 11 \\
B6 & Application des correctifs & 8.8 & ID.RA & 7 \\
B7 & Sensibilisation des agents & 6.3 & PR.AT & 14 \\
B8 & Sécurité des prestataires & 5.19, 5.20 & GV.SC & 15 \\
B9 & Réponse aux incidents & 5.24, 5.26 & RS.MA & 17 \\
B10 & Test d'intrusion annuel & (partiel) & ID.IM & 18 \\
\end{tableaucours}
\vspace{-0.8em}
\begin{center}{\small\itshape Tableau~9.4. La base commune de l'ANPP
(correspondances établies par le cours)}\end{center}

Trois remarques. B5 correspond à deux catégories du CSF~: la sauvegarde
protège les données et sert au rétablissement. B10 n'a pas de mesure
dédiée dans l'ISO/IEC~27002~; le test d'intrusion n'y apparaît que comme un
moyen parmi d'autres, d'où la mention «~partiel~». Enfin, la base ne couvre
pas tout~: plusieurs catégories du CSF (le contexte organisationnel, la
politique, la communication) n'ont pas de contrôle dans cette première
version.

\subsection{Tester une fois, rendre compte trois fois}

\begin{boxlecture}{La base commune de l'ANPP, évaluée en 2027}
\textbf{Constatation.} L'équipe de l'ANPP, appuyée par un cabinet, teste
une fois les dix contrôles~:
\begin{itemize}
    \item \OUI{}~: B1, B5, B9 (exercice de 2026), B10 (rapport de 2027)~;
    \item \NON{}~: B2 (quatre agents payés après leur départ, comptes
          encore actifs~; aucune procédure écrite), B3 (pas de second
          facteur), B7 (aucune sensibilisation depuis 2023)~;
    \item \NEK{}~: B4 (journaux «~aux archives~», non produits), B6
          (aucun relevé des correctifs), B8 (contrats sans clause de
          sécurité, aucun rapport des prestataires).
\end{itemize}
4~\OUI{}, 3~\NON{}, 3~\NEK{}~:
\[
\mu = 0{,}40~;\quad U = 0{,}30~;\quad
H = -(0{,}4\log_2 0{,}4 + 2 \times 0{,}3\log_2 0{,}3)
  \approx 0{,}529 + 0{,}521 + 0{,}521 \approx 1{,}57~\text{bit}.
\]

\textbf{Projection sur le CSF} (par fonction)~: Gouverner 0/0/1 (B8)~;
Identifier 2/0/1 (B1, B10~; B6)~; Protéger 1/3/0 (B5~; B2, B3, B7)~;
Détecter 0/0/1 (B4)~; Répondre 1/0/0 (B9)~; Rétablir 1/0/0 (B5).

\medskip
\textbf{Analyse} (confiance moyenne). La même série de tests donne trois
lectures. Pour le bailleur, le profil actuel est faible en
\emph{Protéger} (des défauts établis) et inconnu en \emph{Gouverner} et
\emph{Détecter}~: le niveau de mise en œuvre est au mieux «~partiel~», et
l'ANPP vise «~reproductible~». Pour la Chambre des comptes, B2 rejoint le
constat des quatre agents payés après leur départ~: le défaut de
désactivation est aussi un défaut de contrôle des paiements. Pour
l'audit réglementaire, les mêmes preuves serviront, référentiel par
référentiel.

\medskip
\hfill\lecture{0{,}40}{0{,}30}{1{,}57}
\end{boxlecture}

\begin{boxalert}
\textbf{Violation typique.} Reporter une cote sur un référentiel par la
seule table, alors que le lien est partiel.

\smallskip
\textbf{Règle~:} un \OUI{} ne se transporte que par un lien complet. Par un
lien partiel, le \OUI{} devient au mieux un \NEK{} sur l'autre référentiel,
jusqu'à ce que la partie non couverte soit testée.
\end{boxalert}

% ============================================================================
\section{Choisir ses Référentiels}
\label{sec:c09-choisir}
% ============================================================================

On ne choisit pas un référentiel pour sa réputation, mais pour ce qu'il
oblige ou permet.

\begin{tableaucours}{L{3.4cm}L{8.8cm}}{%
    \tableauHeader{Critère} &
    \tableauHeader{Question à se poser}
}
Obligation & La loi, le superviseur ou un contrat impose-t-il un
référentiel (PCI DSS, exigence d'un bailleur, audit réglementaire)~? \\
Reconnaissance & Les clients ou partenaires attendent-ils une certification
(ISO/IEC~27001) ou un rapport (SOC~2, ISAE~3402)~? \\
Taille et maturité & L'organisation peut-elle faire vivre un SMSI, ou
doit-elle commencer par une hygiène de base (CIS IG1)~? \\
Langage & Faut-il parler à une direction et à un bailleur (CSF) ou à des
techniciens (SP~800-53, guides de configuration)~? \\
Coût de la preuve & Combien coûte la démonstration de conformité, et
peut-elle être mutualisée (base commune)~? \\
\end{tableaucours}
\vspace{-0.8em}
\begin{center}{\small\itshape Tableau~9.5. Critères de choix d'un référentiel}\end{center}

\begin{boxscenario}{Trois organisations, trois choix}
\textbf{BANQUE SANAGA.} ISO/IEC~27001 certifiée, sur toute la banque~;
PCI DSS pour la monétique~; exigences du superviseur
(Chapitre~\ref{chap:droit})~; rapports ISAE~3402 demandés aux prestataires.

\medskip
\textbf{LiZbiZ.} CIS Controls IG1 d'abord, comme socle vérifiable~;
ISO/IEC~27001 ensuite, si PHARMACAM l'exige au renouvellement suivant.

\medskip
\textbf{ANPP.} Une base commune ISO/IEC~27002, projetée sur le CSF pour le
bailleur et sur les textes pour la Chambre des comptes et l'audit
réglementaire.
\end{boxscenario}

% ============================================================================
\section{Ce que les Référentiels demandent à chaque Métier}
\label{sec:c09-metiers}
% ============================================================================

\begin{boxmetier}{Responsable GRC}
\textbf{Une base, plusieurs rapports}
\begin{enumerate}
    \item Recenser les référentiels imposés et attendus.
    \item Construire une base commune de contrôles, chacun relié aux
          référentiels par des liens qualifiés (complet, partiel).
    \item Tester chaque contrôle une fois, avec des preuves réutilisables.
    \item Produire les rapports par projection, en signalant les liens
          partiels.
\end{enumerate}
\end{boxmetier}

\begin{boxmetier}{Auditeur en cybersécurité}
\textbf{Ne pas auditer la table}

L'auditeur à qui l'on présente une table de correspondance vérifie sa
source, sa version et la nature des liens. Il refuse de reporter une
conformité par un lien partiel, et prélève des contrôles dans les zones
que la base ne couvre pas.
\end{boxmetier}

\begin{boxmetier}{Gestionnaire des risques cyber}
\textbf{Le référentiel ne remplace pas le risque}

Un référentiel donne des mesures~; il ne dit pas lesquelles comptent le
plus pour l'organisation. Le gestionnaire des risques vérifie que les
contrôles de la base traitent les scénarios du registre, et pas seulement
les cases des référentiels.
\end{boxmetier}

\separateur

\begin{boxresume}
\textbf{Ce qu'il faut retenir de ce chapitre~:}
\begin{itemize}
    \item \textbf{NIST CSF~2.0}~: six fonctions, 22 catégories~; un langage,
          des profils et des niveaux, pas d'exigences.
    \item \textbf{Catalogues}~: NIST SP~800-53 (détaillé), CIS Controls (18
          contrôles priorisés, groupe IG1 pour l'hygiène de base), PCI DSS
          (cartes de paiement, contractuel).
    \item \textbf{Rapports d'assurance}~: ISAE~3402 et SOC~2~; seul le
          type~2 dit si les contrôles ont fonctionné.
    \item \textbf{Une correspondance} relie des sujets, pas des
          contenus~; elle a une source, une nature, un sens et une
          version.
    \item \textbf{Une base commune} se teste une fois et se projette sur
          plusieurs référentiels~; un \OUI{} ne se transporte que par un
          lien complet.
\end{itemize}
\end{boxresume}

\begin{boxinfo}
\textbf{Points de contrôle.} Vous avez compris ce chapitre si vous savez~:
\begin{itemize}
    \item placer une pratique dans une fonction du CSF~;
    \item dire ce qu'apporte un rapport de type~2 par rapport au type~1~;
    \item construire une ligne de base commune avec trois référentiels~;
    \item refuser une conformité déduite d'un lien partiel.
\end{itemize}
\end{boxinfo}

% ============================================================================
\section*{Exercices}
% ============================================================================

\begin{boxexo}[title={\ding{45}\quad Exercice~9.1. Quelle fonction du CSF~? \niveaubase}]
Placez chaque pratique dans une fonction du NIST CSF~2.0, et, si possible,
dans une catégorie.
\begin{enumerate}[(a)]
    \item Le conseil fixe l'appétence au risque cyber.
    \item La DSI tient l'inventaire des postes des agences.
    \item Les agents suivent une formation à l'hameçonnage.
    \item Le collecteur de journaux déclenche une alerte sur une connexion
          inhabituelle.
    \item La cellule de crise décide d'isoler les agences touchées.
    \item Les agences rouvrent après restauration des postes.
    \item La banque exige un rapport de type~2 de DIGICORE.
\end{enumerate}
\end{boxexo}

\begin{boxexo}[title={\ding{45}\quad Exercice~9.2. Correspondances partielles \niveaumoyen}]
\begin{enumerate}[(a)]
    \item Pour chacun des contrôles B2, B4 et B10 de l'ANPP, dites si le lien
          avec la mesure ISO/IEC~27002 indiquée est complet ou partiel, et
          pourquoi.
    \item Un consultant affirme que l'ANPP, étant «~conforme à 60~\% de
          l'ISO/IEC~27002~», est «~au niveau~2 du CSF~». Relevez les deux
          erreurs de raisonnement.
    \item B10 est \OUI{}. Quelle cote reporter sur l'ISO/IEC~27002, et
          pourquoi~?
\end{enumerate}
\end{boxexo}

\begin{boxexo}[title={\ding{45}\quad Exercice~9.3. L'hygiène de base de LiZbiZ \niveaumoyen}]
\begin{enumerate}[(a)]
    \item Justifiez le choix du groupe IG1 des CIS Controls pour LiZbiZ, à
          l'aide des critères du Tableau~9.5.
    \item Reliez chacun des constats de l'audit des quatre centres (comptes
          par défaut, services en clair, flux non autorisés, sauvegardes non
          restaurées) à un contrôle CIS.
    \item PHARMACAM demande «~une certification~». Que répondez-vous à
          M.~ETEME~: CIS, ISO/IEC~27001 ou rapport SOC~2~? Justifiez.
\end{enumerate}
\end{boxexo}

\begin{boxexo}[title={\ding{45}\quad Exercice~9.4. Trois rapports pour l'ANPP \niveauexpert}]
À partir des résultats de la base commune (4/3/3)~:
\begin{enumerate}[(a)]
    \item Calculez $\mu$, $U$ et $H$ pour la fonction \emph{Protéger} et
          pour la fonction \emph{Identifier}, et appliquez la règle de
          décision du Chapitre~\ref{chap:conclure}.
    \item Rédigez, en dix lignes, la partie «~profil actuel et profil
          cible~» destinée au bailleur.
    \item Rédigez, en cinq lignes, la note destinée à M.~NJOYA sur le
          contrôle B2, en reliant le défaut de sécurité au défaut de
          paiement.
    \item Proposez trois contrôles à ajouter à la base pour couvrir les
          catégories du CSF qu'elle ignore, avec leur correspondance
          ISO/IEC~27002.
\end{enumerate}
\end{boxexo}

\bigskip

\begin{boxinfo}
\textbf{Pour aller plus loin.}
\begin{itemize}
    \item Le NIST CSF~2.0~\cite{nist2024csf} et ses références informatives.
    \item Les CIS Controls, version~8~\cite{cis2021controls}.
    \item La norme ISAE~3402~\cite{iaasb2011isae3402} et les critères de
          services de confiance de l'AICPA~\cite{aicpa2017tsc}.
\end{itemize}
\end{boxinfo}

\begin{boxinfo}
\textbf{Transition.} Les référentiels sont volontaires ou contractuels. La
loi, elle, s'impose~: quelles organisations doivent être auditées, par
qui, avec quelles garanties pour l'auditeur et l'audité~? Le
Chapitre~\ref{chap:droit} présente le cadre juridique camerounais et la
déontologie de l'auditeur.
\end{boxinfo}
